\section{Mergeable Sketches, Neural Operators, and C-TreePO}
\label{app:v6-mergeable-sketches}

Mergeable summaries are a standard abstraction for answering queries about
large data without keeping the full data in memory.  The data are split into
shards.  Each shard is replaced by a compact state.  The states are then
combined up a tree, and the final query is asked only at the root.  A distinct
count sketch, for example, does not store every item it has seen; it stores a
small state from which the number of distinct items can be estimated.

C-TreePO uses the same state-level pattern for learned compression trees.  The
summarizer/merger \(g\) builds leaf states and merges internal states; the
scorer \(f\) is applied only to the root state.  The central question is
whether the state is a \emph{score-sufficient summary}: after a raw subtree is
replaced by its state, the task score remains recoverable at the root.  In a
classical sketch, the final query plays this role.  In C-TreePO, the root
score plays it: a task score, preference judgment, calibrated reward, or other
task-level value.  This is not a scalar-output claim.  In general, child
answers do not determine the answer at their parent.

\subsection{Mergeable Summaries as Score-Sufficient States}
\label{app:v6-mergeable-pattern}

\citet{Agarwal2013MergeableTODS} write \(S()\) for a summary method.  The
input \(D\) is a dataset or stream multiset: if the same record appears twice,
both appearances count.  The parameter \(\varepsilon\) is the target error
level.  Given \(D\) and \(\varepsilon\),
\[
  S(D,\varepsilon)
\]
denotes a summary state produced for \(D\) at that accuracy.  A state is
\emph{valid} when it makes the three promises required by the sketching
problem:
\begin{enumerate}[leftmargin=1.5em]
\item it is a well-formed state for the summary method;
\item its size is within the advertised bound;
\item the final query answered from the state has error at most
\(\varepsilon\).
\end{enumerate}
The notation is deliberately one-to-many.  Different stream orders or random
choices may produce different valid summaries for the same \(D\).  The size
profile records the largest valid state that can arise on inputs of a given
size:
\[
  k(n,\varepsilon)
  :=
  \max\{\, |S(D,\varepsilon)| : |D|=n,\ S(D,\varepsilon)
    \text{ is valid}\,\},
\]
with \(k(\varepsilon)\) used when the bound is independent of \(n\).

Full mergeability is a state statement.  There is a merge algorithm \(A\) that
takes two already-valid states and returns a valid state for the combined
data:
\[
  A\big(S(D_1,\varepsilon),S(D_2,\varepsilon)\big)
  \quad\text{is a valid}\quad
  S(D_1\uplus D_2,\varepsilon).
\]
Here \(D_1\uplus D_2\) is multiset addition: repeated records remain repeated.
It is the data operation represented by joining two shards or subtrees.  The
final sketch query is applied after the state merge, not at the children.

The distinct-count case gives the basic intuition.  If \(D_1\) and \(D_2\)
are two shards, the two scalar answers ``how many distinct items are in
\(D_1\)?'' and ``how many distinct items are in \(D_2\)?'' do not determine
how many distinct items are in \(D_1\uplus D_2\), because the overlap is
unknown.  A mergeable state keeps enough information about the shard to merge
states first and ask the distinct-count query afterward.  This is why the
state, rather than the child scalar answer, is the object being merged.

To compare this with C-TreePO, unpack the overloaded
\(S(D,\varepsilon)\) notation.  The expression \(S(D,\varepsilon)\) should be
read as ``some valid state produced for \(D\) at accuracy
\(\varepsilon\),'' not as a single canonical object.  Classical sketches
often choose their state format from the target accuracy:
\[
  \mathcal{Z}_\varepsilon,\qquad
  \mathrm{Valid}_\varepsilon,\qquad
  \mathrm{merge}_\varepsilon.
\]
For the rest of the state-level derivation, fix \(\varepsilon\) and suppress
this subscript:
\[
  \mathcal{Z}:=\mathcal{Z}_\varepsilon,\qquad
  \mathrm{Valid}:=\mathrm{Valid}_\varepsilon,\qquad
  \mathrm{merge}:=\mathrm{merge}_\varepsilon.
\]
C-TreePO trees use the same convention.  The value \(\varepsilon\) is the
target root-score error, and \(\mathcal{Z}\) is fixed for the run.

Let \(\mathcal{X}\) be the space of objects that can be scored, such as
strings, token sequences, record multisets, embeddings, or rendered summary
states.  The object \(X\in\mathcal{X}\) is the C-TreePO representation
corresponding to \(D\).  With \(\varepsilon\) fixed, a mergeable summary becomes
\[
  S(D,\varepsilon)
  \quad\leadsto\quad
  z\in\mathcal{Z}
  \ \text{and}\quad
  \mathrm{Valid}(X,z).
\]
The relation \(\mathrm{Valid}(X,z)\) says that \(z\) is an acceptable state
for the represented object \(X\): it is well formed, has the required size,
and preserves the relevant root task value to accuracy \(\varepsilon\).
The mergeability law becomes
\[
  \mathrm{Valid}(X_1,z_1)
  \wedge
  \mathrm{Valid}(X_2,z_2)
  \Longrightarrow
  \mathrm{Valid}
  \big(X_1\uplus X_2,\mathrm{merge}(z_1,z_2)\big).
\]

\begin{table}[t]
\centering
\small
\setlength{\tabcolsep}{4pt}
\begin{tabular}{@{}p{0.28\linewidth}p{0.34\linewidth}p{0.28\linewidth}@{}}
\toprule
Mergeable-summary notation & State-level notation & C-TreePO role \\
\midrule
\(S()\) & summary method / state interface & choice of state space, validity, merge, and query \\
\(S(D,\varepsilon)\) & a state \(z\) with
\(\mathrm{Valid}(X,z)\) & a valid summary state for the represented input \(X\) corresponding to \(D\) \\
\(A(S(D_1,\varepsilon),S(D_2,\varepsilon))\) &
\(\mathrm{merge}(z_1,z_2)\) & internal \(g\)-call on child states \\
\(D_1\uplus D_2\) & represented multiset union \(X_1\uplus X_2\) & data represented by the parent subtree \\
\(k(n,\varepsilon)\) & maximum size of valid states for \(|D|=n\) & state-size profile preserved by the tree \\
final query on \(S(D,\varepsilon)\) &
\(\mathrm{query}(z,\cdot)\) & root score \(f(z_{\mathrm{root}})\) after specializing the query family \\
\bottomrule
\end{tabular}
\caption{Translation from the paper's overloaded summary notation to the
state-level notation used for the C-TreePO correspondence, after fixing
\(\varepsilon\).}
\label{tab:v6-agarwal-translation}
\end{table}

One-way mergeability is the streaming special case: one side of the merge is a
fresh singleton or raw suffix.  Equal-size or restricted variants constrain
which pairs of states may be merged.

For randomized summaries, validity is an event.  A single random seed
\(\omega\in\Omega\) may encode all random choices used by a merge-tree
evaluation.  The paper-level success statement says that, for any merge tree
whose leaves represent \(D\), the root state lies in \(S(D,\varepsilon)\) with
probability at least \(p\).  C-TreePO preserves that probability shape: on the
root-validity event, the valid-state guarantee becomes the desired root score.

For range spaces, the subset notation is literal.  A finite range
space is \((D,\mathcal{R})\), and an \(\varepsilon\)-approximation is a
subset \(S\subseteq D\) satisfying
\[
  \sup_{R\in\mathcal{R}}
  \left|
    \frac{|S\cap R|}{|S|}
    -
    \frac{|D\cap R|}{|D|}
  \right|
  \le \varepsilon.
\]
In the C-TreePO translation, \(X_1\uplus X_2\) is the represented-input union
tracked by the tree, while \(S\) is the selected subset state.  Keeping those
roles separate prevents the mistaken inference that scalar child answers
should themselves merge.

\subsection{C-TreePO Trees as Learned State Machines}
\label{app:v6-ctreepo-state-machine}

A C-TreePO run fixes a tree \(T\), a state space \(\mathcal{Z}\), a state map
\(g\), and a scorer \(f:\mathcal{X}\to Y\).  The state map \(g\) is used in
two places: it summarizes raw leaves, and it merges states that have already
been produced.  Each node \(v\) represents an object
\(X_v\in\mathcal{X}\).  At a leaf, the represented object is the raw block
\(b\), and the tree stores the state
\[
  z_b = g(b).
\]
At an internal node, the represented object is the union of the two child
objects.  The tree does not rescore the children and combine their scalar
answers.  It combines the child states:
\[
  z_v = g(z_L\concat z_R).
\]
Here \(\concat\) denotes the representation supplied to the merger, for
example concatenated text summaries, paired JSON states, stacked embeddings,
or any other fixed representation of the two child states.  After all merges,
the tree prediction is the root score
\[
  \widehat y_T = f(z_{\mathrm{root}}),
\]
and the oracle comparison is against \(f(X_T)\), where \(X_T\) is the object
represented by the full tree.  For learned or LLM-style trees, prompts,
schemas, embeddings, or hidden-state formats are usually fixed before the
audit; \(\varepsilon\) is the certification threshold.  To score a state,
either regard rendered states as lying in \(\mathcal{X}\), so
\(\mathcal{Z}\subseteq\mathcal{X}\), or use a renderer
\(\rho:\mathcal{Z}\to\mathcal{X}\) and interpret \(f(z)\) as \(f(\rho(z))\).

Fix a task distance \(d_Y:Y\times Y\to[0,\infty)\).  Throughout this appendix
\(d_Y\) is taken on the same scale as the summary error guarantee, so a valid
state is score-sufficient when
\[
  \mathrm{Valid}(X,z)
  \Longrightarrow
  d_Y\big(f(z),f(X)\big)\le\varepsilon.
\]
The state need not reconstruct \(X\).  It only has to preserve the score to
which the root scorer is sensitive.  HyperLogLog illustrates the distinction:
the scalar counts \(|A|\) and \(|B|\) do not determine \(|A\cup B|\), but HLL
register states merge by coordinatewise maximum and retain the overlap
information needed for the final distinct-count query
\citep{FlajoletEtAl2007,HeuleEtAl2013}.

Define exact score equivalence by
\[
  u\sim_f v
  \Longleftrightarrow
  d_Y\big(f(u),f(v)\big)=0,
\]
and the \(\varepsilon\)-version by
\[
  u\sim_{f,\varepsilon} v
  \Longleftrightarrow
  d_Y\big(f(u),f(v)\big)\le\varepsilon.
\]
The local laws are the checks that make the tree computation credible.  C1
asks whether a leaf state preserved what matters about its raw block.  C3
asks whether a merged state preserved what matters about the union of the two
children.  C2 asks whether a state that has already been produced can safely
be supplied to \(g\) again without changing its score-relevant content.  Let
\(B_T\) be the composed root score-error bound obtained from the checked
C1/C2/C3 residuals along \(T\).  The tree is certified at accuracy
\(\varepsilon\) when \(B_T\le\varepsilon\).  The exact versions of the laws
are
\begin{align}
\tag{C1}
&\forall\text{ realized leaves }b:\quad g(b)\sim_f b,\\
\tag{C2}
&\forall\text{ produced states }z:\quad g(z)\sim_f z,\\
\tag{C3}
&\forall\text{ sibling nodes }L,R:\quad
  g(z_L\concat z_R)\sim_f X_L\uplus X_R.
\end{align}
In C3, \(z_L\) and \(z_R\) are the produced states for sibling nodes whose
represented objects are \(X_L\) and \(X_R\).  C1 is the valid-leaf condition.
C3 is the valid-merge condition.  C2 is the reuse condition: when an
already-produced state is supplied again to \(g\), it remains in the same
score-relevant validity class.

Table~\ref{tab:v6-local-laws-sketch-link} records the corresponding
mergeable-summary obligations.

\begin{table}[t]
\centering
\scriptsize
\setlength{\tabcolsep}{3pt}
\begin{tabular}{@{}p{0.16\linewidth}p{0.24\linewidth}p{0.25\linewidth}p{0.25\linewidth}@{}}
\toprule
Local law & C-TreePO wording & Mergeable-summary wording & Score-sufficiency reading \\
\midrule
C1: Leaf sufficiency &
The leaf state \(g(b)\) preserves the root score that \(f\) would assign to the
raw block \(b\). &
A state produced for represented input \(X_i\) is valid for \(X_i\). &
The leaf state is score-sufficient for its shard: scoring the state recovers
the shard's task value, exactly or within the induced score tolerance. \\
C3: Merge consistency &
The parent state \(g(z_L\concat z_R)\) preserves the root score of the union
represented by the two children. &
Merging valid states for \(X_1\) and \(X_2\) gives a valid state for
\(X_1\uplus X_2\). &
Score sufficiency is closed under state merge, so a tree of local merges
produces a score-sufficient root state. \\
C2: Re-summary stability &
Re-entering an already-produced state into \(g\) does not change the root
score. &
A valid produced summary may be reused as an input state without leaving the
same validity/score class. &
The state remains score-sufficient under reuse: resummarization stays in the
same score-equivalence class. \\
\bottomrule
\end{tabular}
\caption{C-TreePO local laws as score-sufficiency conditions.  Read C1 as
``the leaf summary kept the score-relevant information,'' C3 as ``the merged
summary kept the score-relevant information for the parent,'' and C2 as
``reusing an already-produced summary is stable.''}
\label{tab:v6-local-laws-sketch-link}
\end{table}

\subsection{The C-TreePO Correspondence}
\label{app:v6-ctreepo-correspondence}

With the state-machine notation in place, the correspondence is direct.  A
classical mergeable summary supplies \(g\): use the summary method at leaves,
use the sketch merge at internal nodes, and score the root state with \(f\).
Conversely, if an exact C-TreePO operator satisfies C1--C3 on its realized
states, then its range carries an induced merge
\[
  \mathcal{Z}_g:=\operatorname{range}(g),\qquad
  z\oplus_g z' := g(z\concat z').
\]
C1 supplies valid leaves, C3 supplies valid parent states under \(\oplus_g\),
and C2 says produced states may be reused without leaving the same
score-relevant class.

\begin{proposition}[State-level nesting]
\label{prop:v6-state-level-correspondence}
After fixing a state representation and score map, state-level mergeable
summaries nest in C-TreePO.  Valid leaves and valid state merges give a valid
root state for every merge tree, and root validity gives the promised
\(\varepsilon\)-score guarantee.
\end{proposition}

\begin{corollary}[Randomized root guarantee]
\label{thm:v6-randomized-root-success}
If a randomized mergeable-summary tree has root-validity probability at
least \(p\), and every valid root state satisfies the promised score
guarantee, then the C-TreePO root score satisfies that guarantee with
probability at least \(p\).  The statement is event-level: no seedwise
validity assumption is required.
\end{corollary}

\begin{remark}[Scalar-query caveat]
\label{thm:v6-scalar-query-caveat}
A scalar child-query merge law is a separate stronger assumption.  Classical
mergeability of states does not imply that child scalar answers determine the
parent scalar answer.  Distinct count is the running example: child scalar
counts omit overlap information, while HLL states retain enough information to
merge.
\end{remark}

This is the precise nesting claim: classical mergeable summaries enter
C-TreePO through relational state-level trees, not through scalar child
answers.

\subsection{Neural Operators as a Representation Route}
\label{app:v6-neural-operator-route}

The neural-operator role is representational.  In the state-machine view, a
neural operator supplies a learned family for \(g\).  Informally, this means
that the summarizer and merger are not hand-written sketch algorithms; they
are parameterized maps trained from examples, oracle feedback, or local-law
losses.  Neural operators are designed to act on structured inputs such as
functions, fields, or high-dimensional states \citep{Kovachki23NeuralOperator}.
If the state space is represented in the operator domain, and if the realized
tree generates a finite collection of leaf, merge, and on-range re-summary
calls, then an appropriate neural-operator class can approximate the ideal
state-building and state-merge maps on that realized call set.

Approximation is not certification.  A flexible architecture may learn a useful
state map, but it is the local-law or projection layer that tells us whether
the learned state behaves like a mergeable summary for the task.  A typical
training objective has the form
\[
  \mathcal{L}
  =
  \mathrm{calibration\_loss}(\widehat f,\fstar)
  +\mathrm{root\_loss}(\widehat f,g)
  +\lambda_1 C1+\lambda_3 C3+\lambda_2 C2.
\]
The coefficients \(\lambda=(\lambda_1,\lambda_2,\lambda_3)\) are fixed
training hyperparameters.  They set how strongly training penalizes local-law
violations; they are not error units.  Certification is the separate statement
that the measured root error is at most the target \(\varepsilon\).
In simulations where exact sketch states are known, one may also add state
imitation losses.  In semantic applications such as the Manifesto task, the
state is not a hand-written register vector, so the local laws and sampled
audit are the observable certification route.

A compact way to state the learned guarantee is this.  Let \(B_T(g,\widehat f)\)
be the root error certified by the checked local laws for the learned scorer
\(\widehat f\), and suppose \(\widehat f\) is uniformly calibrated to the true
oracle \(\fstar\) within \(\varepsilon_f\).  Then
\[
  B_T(g,\widehat f)+2\varepsilon_f\le\varepsilon
  \quad\Longrightarrow\quad
  \text{the learned tree is certified for }\fstar\text{ at }\varepsilon.
\]
The factor \(2\varepsilon_f\) is the two-sided cost of comparing through the
calibrated scorer \(\widehat f\).  A neural-operator architecture therefore
does not by itself define a mergeable summary.  It is covered only when the
realized state machine lies in, or is certified near, the mergeable-state
local-law subspace.

\subsection{Sketch Families and Claim Status}
\label{app:v6-sketch-inventory}

Table~\ref{tab:v6-sketch-inventory} separates theorem-backed claims from
conditional theorem bundles and empirical reference rows.  The distinction is
important for the citation above.  \citet{Agarwal2013MergeableTODS} do not
state an \(\varepsilon\)-kernel theorem, so \(\varepsilon\)-kernel geometry is
adjacent background here unless cited to a separate source.

\begin{table}[t]
\centering
\scriptsize
\setlength{\tabcolsep}{3pt}
\begin{tabular}{@{}p{0.16\linewidth}p{0.26\linewidth}p{0.18\linewidth}p{0.27\linewidth}@{}}
\toprule
Family & Description & Status in this paper & What C-TreePO uses \\
\midrule
HLL & Estimates the number of distinct items without storing the items
themselves; the state is a vector of registers. & theorem-backed + empirical & Register-state summary, pointwise-max merge, and final estimator score; the tree and flat reductions agree when they use the same state representation. \\
Count-Min & Estimates item frequencies with hashed counter tables; the query
asks for the approximate count of one item. & theorem-backed + empirical & Additive table summary and point-frequency score; mergeable but not idempotent. \\
Misra--Gries / SpaceSaving & Find heavy hitters: items whose frequencies are
large enough that they should survive compression. & theorem-backed interfaces & Capacity, mass, and error-envelope arguments for Misra--Gries; SpaceSaving follows by the standard min-counter correspondence. \\
GK quantiles & Maintains a compressed ordered summary so approximate ranks and
quantiles can be answered from a stream. & theorem-backed one-way & Sequential/incremental quantile interface; not claimed as arbitrary full state-state mergeability. \\
KLL quantiles & Randomized quantile sketch that compacts weighted samples and is
designed for mergeable quantile states. & randomized theorem-backed interface & Randomized sized mergeable query-sketch interface with a weighted-mass invariant. \\
Interval/range constructions & Keep small subsets or colored states that approximate many interval or range counts at once. & theorem-backed scaffolds + conditional size-reduction constructions & Same-weight interval concentration, finite VC/Sauer--Shelah trace layer, and low-discrepancy coloring route when supplied by a certificate. \\
Markov and bigram sketches & Store the count information needed to score simple sequence models, such as transition or bigram statistics. & theorem-backed local-law examples & Task-specific score-sufficient states showing the same summarize/merge/score pattern outside classical data-stream queries. \\
CPC, Theta/KMV, Frequent Items, classic quantiles, REQ, t-digest, Tuple, VarOpt & Additional library sketch families used to compare tree reductions against familiar implementations. & empirical-only here & Empirical comparisons only; no theorem is claimed for these additional variants. \\
\(\varepsilon\)-kernel geometry & Geometric approximation family for preserving directional or shape information, separate from the mergeable-summary results of \citet{Agarwal2013MergeableTODS}. & adjacent/background here & Do not attribute to that paper without a separate source. \\
\bottomrule
\end{tabular}
\caption{Mergeable-sketch claim status.  The table records the claims used
here, not the full external literature on each sketch family.}
\label{tab:v6-sketch-inventory}
\end{table}

\subsection{Tree Comparisons Against Classical and Learned Sketches}
\label{app:v6-tree-comparisons}

The empirical comparisons use the correspondence in three ways.

\paragraph{Classical state equality.}
HLL is the tightest check because its state is explicit: a register vector with
pointwise-max merge.  When C-TreePO is instantiated with that state and merge,
the tree-reduced register vector agrees with the flat register vector for every
tested tree shape.  Alternative HLL variants can agree at the estimate level
even when their internal representation changes along different construction
paths.

\begin{figure}[t]
\centering
\includegraphics[width=0.92\linewidth]{classical_sketches_hll_leaf_size.pdf}
\caption{\textbf{Tree reduction, classical HLL, and learned HLL-like states.}
The bottom axis gives leaves per document and the top axis gives the
corresponding token budget per leaf.  \emph{Official/oracle} uses the known HLL
state, merge, and scoring map.  \emph{Learned \(g\) + oracle \(f\)} trains only
the state/merge map while keeping the HLL scoring map fixed.  \emph{Learned
\(f+g\)} jointly learns the latent state, merge, and scoring map.}
\label{fig:v6-hll-parity-curves}
\end{figure}

\paragraph{Learned state with fixed versus learned scoring map.}
The learned HLL rows ask whether a neural \(g\) can recover a classical
state-and-estimator target.  With a fixed classical HLL estimator as \(f\),
merge errors must land in the exact register geometry and can compound with
tree depth.  With learned \(g+f\), the system can choose a different latent
state that preserves the same score.  This is the practical lesson for
semantic C-TreePO tasks.  Exact recovery of a reference sketch imposes a real
scoring constraint; task preservation alone may be better served by a jointly
learned state and scoring map.

\paragraph{Broad sketch suite.}
The broader comparison applies the same tree-reduction protocol to classical
sketches and learned companions across distinct-counting, frequency, quantile,
set, and sampling-style queries.  The table reports representative rows;
larger grids are omitted here to keep the appendix focused.

\begin{figure}[t]
\centering
\includegraphics[width=0.92\linewidth]{classical_sketches_summary.pdf}
\caption{\textbf{Broad sketch leaf-size summary.}
Large-capacity rows are shown by sketch family, with leaves per document and
tokens per leaf shown as paired axes.  The three curves separate the known
classical state and oracle scoring map, a learned state/merge map with known
oracle \(f\), and a jointly learned \(f+g\) model.}
\label{fig:v6-classical-sketches-summary}
\end{figure}

\begin{table}[t]
\centering
\scriptsize
\setlength{\tabcolsep}{4pt}
\begin{tabular}{@{}p{0.22\linewidth}p{0.32\linewidth}p{0.34\linewidth}@{}}
\toprule
Task family & Classical comparison & Learned-tree comparison \\
\midrule
Distinct counting & HLL gives the canonical state-level parity check:
state merge matches flat reduction when the same register representation is
used. & Learned HLL-like states reveal the cost of requiring exact reference
state recovery versus allowing a learned score. \\
Frequency estimation & Count-Min has additive table states and point-query
scores. & Learned companions test whether the tree can preserve frequency
information under repeated merges. \\
Quantiles and ranks & GK and KLL represent the deterministic one-way and
randomized mergeable quantile regimes. & Learned companions test whether a
score-sufficient semantic state can preserve rank information well enough for
root scoring. \\
Set and sampling-style summaries & Additional classical sketches serve as
empirical baselines across set, sampling, and weighted-summary tasks. &
Learned rows are treated as empirical evidence only unless they satisfy the
state-level local laws above. \\
\bottomrule
\end{tabular}
\caption{\textbf{Broad-sketch comparison.}  Classical rows are empirical
references; theorem status is the separate classification in
Table~\ref{tab:v6-sketch-inventory}.}
\label{tab:v6-classical-sketches-compact}
\end{table}

\paragraph{Learning a score-sufficient state from oracle feedback.}
The controlled learned-sketch diagnostic removes semantic ambiguity.  The
oracle asks for a \(k\)-vector of per-type counts, so a state with dimension
\(m<k\) is structurally insufficient.  The learned tree converges when
\(m\ge k\) and plateaus when \(m<k\), matching the hand-designed sketch
boundary.  The Manifesto setting uses the same logic, but the
score-sufficient state is semantic rather than a known count vector.

\begin{figure}[t]
\centering
\includegraphics[width=0.88\linewidth]{learned_sketch_leaf_size_diagnostic.pdf}
\caption{\textbf{Learning a mergeable state from oracle feedback.}
Tree granularity is shown with leaves per document and tokens per leaf.  The
diagnostic compares learned \(g\) with fixed oracle \(f\) and jointly learned \(f+g\) against the
empirical and theoretical HLL floors.}
\label{fig:v6-learned-sketch-curves}
\end{figure}

\subsection{Consequences}
\label{app:v6-what-established}

The preceding discussion establishes the following chain:
\begin{enumerate}[leftmargin=1.5em]
\item Classical mergeable summaries are state-level objects: summarize into
state, merge state, and apply the final query at the root.
\item C-TreePO nests those objects through relational state-level trees, not
through scalar child-answer merging.
\item Randomized mergeable summaries nest with the same probability as their
root-validity event.
\item Exact local laws induce a mergeable state scheme on the realized range
of \(g\), and known sketches supply witnesses for those laws directly.
\item Neural operators provide a representation family for the state-building
and state-merge maps; the C-TreePO local laws and audit are the
certification layer.
\item The empirical suite checks both sides: exact collapse to classical HLL
when the state is fixed, and learned recovery/structural failure when the
state is trained from oracle feedback.
\end{enumerate}

This is the sense in which the mergeable-sketch subspace described by
C-TreePO captures the classical mergeable-summary setting while extending it to
learned semantic states such as the Manifesto application in
Section~\ref{sec:min-manifesto-results}.
